% ========================
% PACOTES
% ========================

\usepackage{graphicx}
\usepackage{amsmath}
\usepackage{xcolor}
\usepackage[a4paper, margin=3cm]{geometry}
\usepackage[colorlinks=true, linkcolor=blue]{hyperref}
\usepackage{bookmark}
\usepackage{listings}
\usepackage{float}
\usepackage{booktabs}
\usepackage{amsfonts}

% ========================
% CONFIGURAÇÕES GERAIS
% ========================

\setlength{\parindent}{0pt}
\setlength{\parskip}{2pt}

% ========================
% CORES DO CÓDIGO
% ========================

\definecolor{codebg}{RGB}{245,245,245}
\definecolor{codegray}{RGB}{100,100,100}
\definecolor{codeblue}{RGB}{0,0,180}
\definecolor{codegreen}{RGB}{0,120,0}

% ========================
% ESTILO PYTHON
% ========================

\lstdefinestyle{mypython}{
    language=Python,
    backgroundcolor=\color{codebg},
    basicstyle=\ttfamily\small,
    keywordstyle=\color{codeblue}\bfseries,
    commentstyle=\color{codegreen}\itshape,
    stringstyle=\color{red},
    numbers=left,
    numberstyle=\tiny\color{codegray},
    stepnumber=1,
    numbersep=8pt,
    frame=single,
    breaklines=true,
    showstringspaces=false,
    tabsize=4
}

\lstnewenvironment{pythoncode}
{
  \lstset{style=mypython}
}
{}